% VAL-013 -- generado por herramientas/gen_casos.py; editar el texto libremente.
\section{VAL-013: Modal: 1 GDL equivalente (M1)}\label{sec:VAL-013}
\begin{ficha}
Fuente & [S] catálogo MAM-253 y \cite{chopra}\\
Ticket & MAM-103\\
Engine & commit \texttt{\EngineCommit}\\
Binario & SHA-256 \texttt{\EngineBinSHA}\\
Entrada & SHA-256 \texttt{\EntradaSHA{VAL-013}}\\
Tolerancia & $|c-r|\le A+R|r|$, $A=0$, $R=0.001$ (clase consolidada; orden de ejecución DESCONOCIDO; el marco también distingue preregistro, DESCONOCIDO y VAL-037 tras la r1 de \#23)\\
Estado de datos & \directlua{mamut.estado('VAL-013')}\\
\end{ficha}

\paragraph{Modelo.}
\begin{center}\diag{sdof}\end{center}

\paragraph{Referencia y derivación.} Derivación propia pendiente de redactar a partir del fixture del caso (provenance del Engine). % DERIV-PENDIENTE


\paragraph{Reproducción.}
\reproducir{VAL-013}

% Q2-PLACEHOLDER INICIO (VAL-013)
% Solo Q2 escribe datos/VAL-013/opensees.csv y el texto de modelo OpenSees de este bloque,
% en su propio PR. J2 no lee ni escribe valores de OpenSees.
\paragraph{Modelo de contraste (OpenSees).} \QdosPendiente
% Q2-PLACEHOLDER FIN (VAL-013)

\paragraph{Resultados.}
\begin{center}
\begin{tikzpicture}
\begin{axis}[grafico, xlabel={modo}, ylabel={$T$ [s]}]
\directlua{mamut.serie('VAL-013','referencia','x','y','ref','Referencia analítica')}
\directlua{mamut.serie('VAL-013','engine','x','y','eng','MAMUT Engine')}
\directlua{mamut.serie('VAL-013','opensees','x','y','ops','OpenSees (contraste)')}
\end{axis}
\end{tikzpicture}
\end{center}
\begin{center}\small
\begin{tabular}{rrrrrrl}\toprule
modo & Ref. & Engine & OpenSees & Error rel. & Tol. & Veredicto\\\midrule
\directlua{mamut.tabla('VAL-013','x','y',0,0.001)}
\bottomrule\end{tabular}
\end{center}
